\subsection{Теорема Гильберта--Шмидта}

\begin{definition}[Собственное подпространство]
	Пусть \(H\) --- комплексное гильбертово пространство,
	\[
	A\in\mathcal{L}(H),
	\qquad
	\lambda\in\mathbb{C}.
	\]
	Собственным подпространством, соответствующим \(\lambda\), называется
	\[
	H_\lambda=\ker(A-\lambda I).
	\]
	Если \(x\in H_\lambda\), \(x\neq0\), то
	\[
	Ax=\lambda x.
	\]
\end{definition}

\begin{lemma}[Наличие ненулевого собственного значения]
	Пусть \(H\) --- комплексное гильбертово пространство,
	\[
	A\in\mathcal{K}(H)
	\]
	--- компактный самосопряженный оператор. Если
	\[
	A\neq0,
	\]
	то у \(A\) существует ненулевое собственное значение.
\end{lemma}

\begin{proof}
	Так как \(A\) самосопряжен, по теореме о спектре самосопряженного оператора
	\[
	r(A)=\|A\|.
	\]
	Поскольку
	\[
	A\neq0,
	\]
	имеем
	\[
	\|A\|\neq0.
	\]
	Значит,
	\[
	r(A)>0.
	\]
	
	По той же теореме
	\[
	r(A)=\max\{|m_-|,|m_+|\},
	\]
	причём
	\[
	m_-,m_+\in\sigma(A).
	\]
	Следовательно, хотя бы одно из чисел \(m_-\), \(m_+\) отлично от нуля. Обозначим его через
	\[
	\lambda.
	\]
	Тогда
	\[
	\lambda\in\sigma(A),
	\qquad
	\lambda\neq0.
	\]
	
	По теореме о компактных самосопряженных операторах всякая ненулевая точка спектра является
	собственным значением. Поэтому \(\lambda\) --- ненулевое собственное значение оператора \(A\).
\end{proof}

\begin{theorem}[Гильберта--Шмидта]
	Пусть \(H\) --- сепарабельное комплексное гильбертово пространство, и
	\[
	A\in\mathcal{K}(H)
	\]
	--- компактный самосопряженный оператор. Тогда в \(H\) существует ортонормированный базис,
	состоящий из собственных векторов оператора \(A\).
\end{theorem}

\begin{proof}
	Докажем, что собственных векторов оператора \(A\) достаточно, чтобы породить всё пространство.
	
	\textbf{1. Отдельно рассмотрим тривиальный случай.}
	
	Если
	\[
	A=0,
	\]
	то любой ненулевой вектор является собственным вектором с собственным значением \(0\).
	Так как \(H\) сепарабельно, в \(H\) существует ортонормированный базис. Он автоматически
	состоит из собственных векторов оператора \(A\).
	
	Далее считаем, что
	\[
	A\neq0.
	\]
	
	\textbf{2. Соберём собственные векторы, отвечающие ненулевым собственным значениям.}
	
	По свойствам собственных значений компактного оператора ненулевых собственных значений
	не более чем счётно. Обозначим их через
	\[
	\lambda_1,\lambda_2,\ldots
	\]
	с учётом возможной конечности этого набора.
	
	Для каждого \(\lambda_j\neq0\) собственное подпространство
	\[
	H_{\lambda_j}=\ker(A-\lambda_j I)
	\]
	конечномерно.
	
	\begin{itemize}
		\item \textbf{Выберем базисы в собственных подпространствах.}
		
		В каждом \(H_{\lambda_j}\) выберем ортонормированный базис:
		\[
		e^{(j)}_1,\ldots,e^{(j)}_{m_j}.
		\]
		
		\item \textbf{Объединим эти базисы.}
		
		Все выбранные векторы являются собственными векторами оператора \(A\). Кроме того,
		собственные векторы, отвечающие различным собственным значениям самосопряженного
		оператора, ортогональны.
		
		Значит, система всех выбранных векторов ортонормирована.
	\end{itemize}
	
	Обозначим через \(M\) замкнутое подпространство, натянутое на все эти векторы:
	\[
	M=
	\overline{\operatorname{span}}
	\{e^{(j)}_k:\lambda_j\neq0,\ 1\leqslant k\leqslant m_j\}.
	\]
	
	\textbf{3. Покажем, что на \(M^\perp\) оператор \(A\) равен нулю.}
	
	Сначала заметим, что \(M\) инвариантно относительно \(A\). Действительно, на каждом собственном
	векторе
	\[
	Ae^{(j)}_k=\lambda_j e^{(j)}_k.
	\]
	Значит,
	\[
	A(M)\subset M.
	\]
	
	Так как \(A\) самосопряжен, из инвариантности \(M\) следует инвариантность \(M^\perp\):
	\[
	A(M^\perp)\subset M^\perp.
	\]
	
	Рассмотрим сужение
	\[
	\widetilde A=A|_{M^\perp}:M^\perp\to M^\perp.
	\]
	Оператор \(\widetilde A\) компактен и самосопряжен.
	
	\begin{itemize}
		\item \textbf{Предположим, что \(\widetilde A\neq0\).}
		
		Тогда по лемме у \(\widetilde A\) существует ненулевое собственное значение \(\mu\neq0\).
		Значит, существует ненулевой вектор
		\[
		u\in M^\perp,
		\]
		такой что
		\[
		\widetilde A u=\mu u.
		\]
		
		\item \textbf{Получим противоречие.}
		
		Так как \(\widetilde A\) --- это сужение \(A\), имеем
		\[
		Au=\mu u.
		\]
		Значит, \(u\) --- собственный вектор оператора \(A\), отвечающий ненулевому собственному
		значению \(\mu\).
		
		По построению все собственные векторы, отвечающие ненулевым собственным значениям,
		лежат в \(M\). Следовательно,
		\[
		u\in M.
		\]
		Но одновременно
		\[
		u\in M^\perp.
		\]
		Значит,
		\[
		u\in M\cap M^\perp=\{0\},
		\]
		что противоречит выбору \(u\neq0\).
	\end{itemize}
	
	Следовательно,
	\[
	\widetilde A=0.
	\]
	То есть
	\[
	Ax=0
	\qquad
	\forall x\in M^\perp.
	\]
	Значит,
	\[
	M^\perp\subset\ker A.
	\]
	
	\textbf{4. Добавим собственные векторы для собственного значения \(0\).}
	
	С другой стороны, если
	\[
	x\in\ker A,
	\]
	то
	\[
	Ax=0.
	\]
	Для любого собственного вектора \(e\), отвечающего ненулевому собственному значению \(\lambda\),
	по ортогональности собственных векторов самосопряженного оператора получаем
	\[
	(x,e)=0.
	\]
	Значит,
	\[
	\ker A\subset M^\perp.
	\]
	Итак,
	\[
	M^\perp=\ker A.
	\]
	
	Так как \(\ker A\) --- замкнутое подпространство сепарабельного гильбертова пространства,
	в нём существует ортонормированный базис. Обозначим его через
	\[
	f_1,f_2,\ldots
	\]
	с учётом возможной конечности или пустоты этого набора.
	
	Все эти векторы являются собственными векторами оператора \(A\), соответствующими собственному
	значению \(0\).
	
	\textbf{5. Получим ортонормированный базис всего пространства.}
	
	Мы получили ортогональное разложение
	\[
	H=M\oplus M^\perp
	=
	M\oplus\ker A.
	\]
	В \(M\) уже выбран ортонормированный базис из собственных векторов, отвечающих ненулевым
	собственным значениям. В \(\ker A\) выбран ортонормированный базис из собственных векторов,
	отвечающих собственному значению \(0\).
	
	Объединяя эти две системы, получаем ортонормированную систему
	\[
	\{e^{(j)}_k\}\cup\{f_s\}.
	\]
	Она полна в \(H\), потому что порождает
	\[
	M\oplus\ker A=H.
	\]
	Следовательно, это ортонормированный базис \(H\), состоящий из собственных векторов оператора
	\(A\).
\end{proof}